\section{Supplement on sheaves}
\label{section:0.3}

\subsection{Sheaves with values in a category}
\label{subsection:0.3.1}

\begin{env}[3.1.1]
\label{0.3.1.1}
Let $\mathbf{K}$ be a category, $(A_\alpha)_{\alpha\in I}$,
$(A_{\alpha\beta})_{(\alpha,\beta)\in I\times I}$ two families of objects of
$\mathbf{K}$ such that $A_{\beta\alpha}=A_{\alpha\beta}$, and
$(\rho_{\alpha\beta})_{(\alpha,\beta)\in I\times I}$ a family of morphisms
$\rho_{\alpha\beta}:A_\alpha\to A_{\alpha\beta}$. We shall say that a pair
consisting of an object $A$ of $\mathbf{K}$ and a family of morphisms
$\rho_\alpha:A\to A_\alpha$ is a \emph{solution of the universal problem} defined
by the families $(A_\alpha)$, $(A_{\alpha\beta})$ and
$(\rho_{\alpha\beta})$ if, for every object $B$ of $\mathbf{K}$, the map
which sends every $f\in\operatorname{Hom}(B,A)$ to the family
$(\rho_\alpha\circ f)\in\prod_\alpha\operatorname{Hom}(B,A_\alpha)$ is a
\emph{bijection} of $\operatorname{Hom}(B,A)$ onto the set of the $(f_\alpha)$
such that
$\rho_{\alpha\beta}\circ f_\alpha=\rho_{\beta\alpha}\circ f_\beta$ for every
pair of indices $(\alpha,\beta)$. One sees at once that if such a
solution exists, it is unique up to isomorphism.
\end{env}

\begin{env}[3.1.2]
\label{0.3.1.2}
We shall not recall the definition of a \emph{presheaf}
$U\to\mathcal{F}(U)$ on a topological space $X$ with values in a
category $\mathbf{K}$ (G, I, 1.9); we shall say that such a presheaf is a
\emph{sheaf with values in $\mathbf{K}$} if it satisfies the following
axiom:

\smallskip
\noindent
\textup{(F)} \emph{For every covering $(U_\alpha)$ of an open subset $U$ of $X$
by open subsets $U_\alpha$ contained in $U$, if we denote by $\rho_\alpha$
(resp. $\rho_{\alpha\beta}$) the restriction morphism}
\[
  \mathcal{F}(U)\to\mathcal{F}(U_\alpha)
  \qquad
  \bigl(\text{resp. }\mathcal{F}(U_\alpha)
  \to\mathcal{F}(U_\alpha\cap U_\beta)\bigr),
\]

\oldpage[0\textsubscript{I}]{24}

\emph{the pair consisting of $\mathcal{F}(U)$ and the family $(\rho_\alpha)$ is a
solution of the universal problem for $(\mathcal{F}(U_\alpha))$,
$(\mathcal{F}(U_\alpha\cap U_\beta))$ and $(\rho_{\alpha\beta})$ (3.1.1)}%
\footnote{This is a special case of the general notion of (non-filtered)
\emph{projective limit} (\emph{see} (T, I, 1.8) and the book in preparation
announced in the Introduction).}.

It amounts to the same to say that, for every object $T$ of $\mathbf{K}$, the
family $U\to\operatorname{Hom}(T,\mathcal{F}(U))$ is a \emph{sheaf
of sets}.
\end{env}

\begin{env}[3.1.3]
\label{0.3.1.3}
Suppose that $\mathbf{K}$ is the category defined by a ``species of
structure with morphisms'' $\Sigma$, the objects of $\mathbf{K}$ thus being
the sets endowed with structures of species $\Sigma$ and the morphisms those of
$\Sigma$. Suppose moreover that the category $\mathbf{K}$ satisfies the
following condition:

\smallskip
\noindent
\textup{(E)} If $(A,(\rho_\alpha))$ is a solution of a universal mapping
problem \emph{in the category $\mathbf{K}$} for families
$(A_\alpha)$, $(A_{\alpha\beta})$, $(\rho_{\alpha\beta})$, then it is also
a solution of the universal mapping problem for the same families
\emph{in the category of sets} (that is, when $A$,
the $A_\alpha$ and the $A_{\alpha\beta}$ are regarded as sets, and the $\rho_\alpha$ and
$\rho_{\alpha\beta}$ as maps)%
\footnote{One can prove that this also means that the canonical functor
$\mathbf{K}\to(\mathrm{Ens})$ \emph{commutes with projective limits} (not
necessarily filtered).}.

Under these conditions, condition (F) implies that, regarded as a
presheaf \emph{of sets}, $U\to\mathcal{F}(U)$ is a \emph{sheaf}.
Moreover, for a map $u:T\to\mathcal{F}(U)$ to be a morphism of
$\mathbf{K}$, it is necessary and sufficient, by (F), that each map
$\rho_\alpha\circ u$ be a morphism $T\to\mathcal{F}(U_\alpha)$, which
means that the structure of species $\Sigma$ on $\mathcal{F}(U)$ is the
\emph{initial structure} for the morphisms $\rho_\alpha$. Conversely,
suppose that a presheaf $U\to\mathcal{F}(U)$ on $X$ with values in
$\mathbf{K}$ is a \emph{sheaf of sets} and satisfies the preceding
condition; it is then clear that it satisfies (F), and hence is a
\emph{sheaf with values in $\mathbf{K}$}.
\end{env}

\begin{env}[3.1.4]
\label{0.3.1.4}
When $\Sigma$ is the species of structure of group or of ring, the fact that
the presheaf $U\to\mathcal{F}(U)$ with values in $\mathbf{K}$ is a
\emph{sheaf of sets} implies \emph{ipso facto} that it is a
\emph{sheaf with values in $\mathbf{K}$} (in other words, a sheaf of
groups or of rings in the sense of (G))%
\footnote{This is because in the category $\mathbf{K}$, every morphism
which is a \emph{bijection} (as a map of sets) is an
\emph{isomorphism}. This is no longer true when $\mathbf{K}$ is, for example, the
category of topological spaces.}. But this is no longer the
case when, for example, $\mathbf{K}$ is the category of \emph{topological
rings} (with continuous homomorphisms as morphisms): a
sheaf with values in $\mathbf{K}$ is a sheaf of rings
$U\to\mathcal{F}(U)$ such that, for every open subset $U$ and every covering of $U$
by open subsets $U_\alpha\subset U$, the topology of the ring $\mathcal{F}(U)$
is the \emph{coarsest} making the homomorphisms
$\mathcal{F}(U)\to\mathcal{F}(U_\alpha)$ continuous. In this case we shall say that
$U\to\mathcal{F}(U)$, regarded as a sheaf of rings (without topology),
\emph{underlies} the sheaf of topological rings
$U\to\mathcal{F}(U)$. The morphisms
$u_V:\mathcal{F}(V)\to\mathcal{G}(V)$ ($V$ an arbitrary open subset of $X$) of
sheaves of topological rings are therefore homomorphisms of the underlying sheaves
of rings such that $u_V$ is \emph{continuous} for every open subset
$V\subset X$; to distinguish them from arbitrary homomorphisms of the
underlying sheaves of rings, we shall call them \emph{continuous}
homomorphisms of sheaves of topological rings. There are analogous definitions and
conventions for sheaves of topological spaces or of topological
groups.
\end{env}

\oldpage[0\textsubscript{I}]{25}

\begin{env}[3.1.5]
\label{0.3.1.5}
It is clear that for every category $\mathbf{K}$, if $\mathcal{F}$ is a
presheaf (resp. a sheaf) on $X$ with values in $\mathbf{K}$ and $U$ an
open subset of $X$, the $\mathcal{F}(V)$ for the open subsets $V\subset U$ constitute
a presheaf (resp. a sheaf) with values in $\mathbf{K}$, which is
called the presheaf (resp. sheaf) \emph{induced} by $\mathcal{F}$ on $U$ and
which is denoted $\mathcal{F}|U$.

For every morphism $u:\mathcal{F}\to\mathcal{G}$ of presheaves on $X$ with
values in $\mathbf{K}$, we shall denote by $u|U$ the morphism
$\mathcal{F}|U\to\mathcal{G}|U$ formed by the $u_V$ for $V\subset U$.
\end{env}

\begin{env}[3.1.6]
\label{0.3.1.6}
Suppose now that the category $\mathbf{K}$ admits \emph{inductive
limits} (T, 1.8); then, for every presheaf (and in particular every
sheaf) $\mathcal{F}$ on $X$ with values in $\mathbf{K}$ and every $x\in X$,
one can define the \emph{stalk} $\mathcal{F}_x$ as the object of $\mathbf{K}$
which is the inductive limit of the $\mathcal{F}(U)$ over the set, filtered (for
$\supset$), of open neighbourhoods $U$ of $x$ in $X$, with respect to the morphisms
$\rho_U^V:\mathcal{F}(V)\to\mathcal{F}(U)$. If
$u:\mathcal{F}\to\mathcal{G}$ is a morphism of presheaves with values in
$\mathbf{K}$, one defines, for every $x\in X$, the morphism
$u_x:\mathcal{F}_x\to\mathcal{G}_x$ as the inductive limit of the
$u_U:\mathcal{F}(U)\to\mathcal{G}(U)$ over the set of open neighbourhoods
of $x$; one thus defines $\mathcal{F}_x$ as a covariant functor in
$\mathcal{F}$, with values in $\mathbf{K}$, for every $x\in X$.

When $\mathbf{K}$ is moreover defined by a species of structure with
morphisms $\Sigma$, the elements of $\mathcal{F}(U)$ are also called the
\emph{sections over $U$} of a \emph{sheaf $\mathcal{F}$} with values in
$\mathbf{K}$, and one then writes $\Gamma(U,\mathcal{F})$ instead of
$\mathcal{F}(U)$; for $s\in\Gamma(U,\mathcal{F})$ and $V$ an open subset contained in
$U$, one writes $s|V$ instead of $\rho_V^U(s)$; for every $x\in U$, the
canonical image of $s$ in $\mathcal{F}_x$ is the \emph{germ} of $s$ at the point $x$,
denoted $s_x$ (\emph{we shall never use the notation $s(x)$ in this sense},
this notation being reserved for another notion relating to the
particular sheaves which will be considered in this Treatise (5.5.1)).

If then $u:\mathcal{F}\to\mathcal{G}$ is a morphism of sheaves with values
in $\mathbf{K}$, we shall write $u(s)$ instead of $u_V(s)$ for every
$s\in\Gamma(V,\mathcal{F})$.

If $\mathcal{F}$ is a sheaf of commutative groups, or of rings, or of
modules, the set of $x\in X$ such that
$\mathcal{F}_x\neq\{0\}$ is called the \emph{support} of $\mathcal{F}$, denoted
$\operatorname{Supp}(\mathcal{F})$; this set is not necessarily
closed in $X$.

When $\mathbf{K}$ is defined by a species of structure with morphisms,
\emph{we shall systematically refrain from bringing in the point of
view of ``étalé spaces''} as regards sheaves with values in
$\mathbf{K}$; in other words, we shall never regard a sheaf as
a topological space (nor even as the set which is the union of its stalks), and
neither shall we regard a morphism
$u:\mathcal{F}\to\mathcal{G}$ of such sheaves on $X$ as a continuous
map of topological spaces.
\end{env}

\subsection{Presheaves on a basis of open subsets}
\label{subsection:0.3.2}

\begin{env}[3.2.1]
\label{0.3.2.1}
In what follows we shall restrict ourselves to categories $\mathbf{K}$
admitting \emph{projective limits} (generalized ones, that is,
corresponding to preordered sets which are not necessarily filtered,
cf.\ (T, 1.8)). Let $X$ be a topological space and $\mathfrak{B}$ a basis
of open subsets for the topology of $X$. We shall call a \emph{presheaf on
$\mathfrak{B}$ with values in $\mathbf{K}$} a family of objects
$\mathcal{F}(U)\in\mathbf{K}$, attached to each $U\in\mathfrak{B}$, together with a
family of morphisms
$\rho_U^V:\mathcal{F}(V)\to\mathcal{F}(U)$ defined for every pair $(U,V)$
of elements of $\mathfrak{B}$ such that $U\subset V$,

\oldpage[0\textsubscript{I}]{26}

subject to the conditions $\rho_U^U=\text{identity}$ and
$\rho_U^W=\rho_U^V\circ\rho_V^W$ if $U$, $V$, $W$ in $\mathfrak{B}$ are
such that $U\subset V\subset W$. One can associate with it a \emph{presheaf with
values in $\mathbf{K}$}: $U\to\mathcal{F}'(U)$ in the ordinary sense, by
taking, for every open subset $U$,
\[
  \mathcal{F}'(U)=\varprojlim_{V\subset U}\mathcal{F}(V),
\]
where $V$ runs through the ordered set (for $\subset$, \emph{not filtered in
general}) of the sets $V\in\mathfrak{B}$ such that $V\subset U$; for the
$\mathcal{F}(V)$ form a projective system for the $\rho_V^W$
($V\subset W\subset U$, $V\in\mathfrak{B}$, $W\in\mathfrak{B}$). Indeed,
if $U$, $U'$ are two open subsets of $X$ such that $U\subset U'$, one defines
$\rho'_U{}^{U'}$ as the projective limit (for $V\subset U$) of the canonical morphisms
$\mathcal{F}'(U')\to\mathcal{F}(V)$, in other words the unique
morphism $\mathcal{F}'(U')\to\mathcal{F}'(U)$ which, composed with the
canonical morphisms $\mathcal{F}'(U)\to\mathcal{F}(V)$, gives the canonical
morphisms $\mathcal{F}'(U')\to\mathcal{F}(V)$; the verification of the
transitivity of $\rho'_U{}^{U'}$ is then immediate. Moreover, if
$U\in\mathfrak{B}$, the canonical morphism
$\mathcal{F}'(U)\to\mathcal{F}(U)$ is an isomorphism, which allows us to identify
these two objects%
\footnote{If $X$ is a \emph{Noetherian} space, one can still define
$\mathcal{F}'(U)$ and show that it is a presheaf (in the ordinary sense)
when one assumes only that $\mathbf{K}$ admits projective limits
for \emph{finite} projective systems. Indeed, if $U$ is an arbitrary open subset
of $X$, there is a \emph{finite} covering $(V_i)$ of $U$ consisting
of sets of $\mathfrak{B}$; for every pair $(i,j)$ of indices, let
$(V_{ijk})$ be a finite covering of $V_i\cap V_j$ consisting of sets of
$\mathfrak{B}$. Let $I$ be the set consisting of the $i$ and of the triples $(i,j,k)$,
ordered by the relations $i>(i,j,k)$, $j>(i,j,k)$ only; one then takes
for $\mathcal{F}'(U)$ the projective limit of the system of the $\mathcal{F}(V_i)$
and $\mathcal{F}(V_{ijk})$; one easily checks that this does not depend on the
coverings $(V_i)$ and $(V_{ijk})$ and that $U\to\mathcal{F}'(U)$ is a
presheaf.}.
\end{env}

\begin{env}[3.2.2]
\label{0.3.2.2}
For the presheaf $\mathcal{F}'$ thus defined to be a \emph{sheaf},
it is necessary and sufficient that the presheaf $\mathcal{F}$ on $\mathfrak{B}$
satisfy the condition:

\smallskip
\noindent
\textup{(F\textsubscript{0})} \emph{For every covering $(U_\alpha)$ of
$U\in\mathfrak{B}$ by sets $U_\alpha\in\mathfrak{B}$ contained in
$U$, and for every object $T\in\mathbf{K}$, the map which sends every
$f\in\operatorname{Hom}(T,\mathcal{F}(U))$ to the family
\[
  (\rho_{U_\alpha}^{U}\circ f)
  \in\prod_\alpha\operatorname{Hom}(T,\mathcal{F}(U_\alpha))
\]
is a bijection of $\operatorname{Hom}(T,\mathcal{F}(U))$ onto the set of the
$(f_\alpha)$ such that
$\rho_V^{U_\alpha}\circ f_\alpha=\rho_V^{U_\beta}\circ f_\beta$ for every
pair of indices $(\alpha,\beta)$ and every $V\in\mathfrak{B}$ such that
$V\subset U_\alpha\cap U_\beta$}%
\footnote{This also means that the pair consisting of $\mathcal{F}(U)$ and the
$\rho_\alpha=\rho_{U_\alpha}^{U}$ is a \emph{solution of the universal problem}
defined (3.1.1) by the data
$A_\alpha=\mathcal{F}(U_\alpha)$,
$A_{\alpha\beta}=\prod\mathcal{F}(V)$ (for the $V\in\mathfrak{B}$ such that
$V\subset U_\alpha\cap U_\beta$) and
$\rho_{\alpha\beta}=(\rho_V''):\mathcal{F}(U_\alpha)\to
\prod\mathcal{F}(V)$ defined by the condition that for
$V\in\mathfrak{B}$, $V'\in\mathfrak{B}$, $W\in\mathfrak{B}$,
$V\cup V'\subset U_\alpha\cap U_\beta$, $W\subset V\cap V'$,
$\rho_W^V\circ\rho_V''=\rho_W^{V'}\circ\rho_{V'}''$.}.

The condition is obviously necessary. To show that it is sufficient,
let us first consider a second basis $\mathfrak{B}'$ of the topology of $X$,
\emph{contained in $\mathfrak{B}$}, and show that if $\mathcal{F}''$
denotes the presheaf deduced from the subfamily
$(\mathcal{F}(V))_{V\in\mathfrak{B}'}$, then $\mathcal{F}''$ is
\emph{canonically isomorphic} to $\mathcal{F}'$. Indeed, first of all, the
projective limit (for $V\in\mathfrak{B}'$, $V\subset U$) of the canonical
morphisms $\mathcal{F}'(U)\to\mathcal{F}(V)$ is a morphism
$\mathcal{F}'(U)\to\mathcal{F}''(U)$ for every open subset $U$. If
$U\in\mathfrak{B}$, this morphism is an isomorphism, since by hypothesis the
canonical morphisms $\mathcal{F}''(U)\to\mathcal{F}(V)$ for
$V\in\mathfrak{B}'$, $V\subset U$, factorize as
\[
  \mathcal{F}''(U)\to\mathcal{F}(U)\to\mathcal{F}(V),
\]
and it is immediate that the composites of the morphisms
$\mathcal{F}(U)\to\mathcal{F}''(U)$ and
$\mathcal{F}''(U)\to\mathcal{F}(U)$ thus defined are the identities. This
being so, for every open subset $U$, the morphisms
$\mathcal{F}''(U)\to\mathcal{F}''(W)=\mathcal{F}(W)$ for
$W\in\mathfrak{B}$ and $W\subset U$ satisfy the conditions characterizing the
projective limit of the $\mathcal{F}(W)$ ($W\in\mathfrak{B}$, $W\subset U$), which
proves our assertion, taking into account the uniqueness of a projective limit
up to isomorphism.

That said, let $U$ be an arbitrary open subset of $X$, $(U_\alpha)$ a covering
of $U$ by open subsets contained in $U$, and let $\mathfrak{B}'$ be the
subfamily of $\mathfrak{B}$ consisting of the sets

\oldpage[0\textsubscript{I}]{27}

of $\mathfrak{B}$ contained in at least one $U_\alpha$; it is clear that
$\mathfrak{B}'$ is again a basis of the topology of $U$, so
$\mathcal{F}'(U)$ (resp. $\mathcal{F}''(U_\alpha)$) is the projective limit of the
$\mathcal{F}(V)$ for $V\in\mathfrak{B}'$ and $V\subset U$ (resp.
$V\subset U_\alpha$); axiom (F) is then verified at once by the
definition of the projective limit.

When (F\textsubscript{0}) is satisfied, we shall say, by abuse of language, that
the presheaf $\mathcal{F}$ on the basis $\mathfrak{B}$ is a \emph{sheaf}.
\end{env}

\begin{env}[3.2.3]
\label{0.3.2.3}
Let $\mathcal{F}$, $\mathcal{G}$ be two presheaves on the basis
$\mathfrak{B}$ with values in $\mathbf{K}$; a \emph{morphism}
$u:\mathcal{F}\to\mathcal{G}$ is defined as a family
$(u_V)_{V\in\mathfrak{B}}$ of morphisms
$u_V:\mathcal{F}(V)\to\mathcal{G}(V)$ satisfying the usual compatibility
conditions with the restriction morphisms $\rho_V^W$. With
the notation of (3.2.1), one deduces from it a morphism
$u':\mathcal{F}'\to\mathcal{G}'$ of the corresponding (ordinary) presheaves
by taking for $u_U'$ the projective limit of the $u_V$ for
$V\in\mathfrak{B}$ and $V\subset U$; the verification of the
compatibility conditions with $\rho'_U{}^{U'}$ follows from the functorial character of the
projective limit.
\end{env}

\begin{env}[3.2.4]
\label{0.3.2.4}
If the category $\mathbf{K}$ admits inductive limits, and if $\mathcal{F}$
is a presheaf on the basis $\mathfrak{B}$ with values in $\mathbf{K}$,
then for every $x\in X$ the neighbourhoods of $x$ belonging to $\mathfrak{B}$ form
a cofinal set (for $\supset$) in the set of neighbourhoods of $x$;
hence, if $\mathcal{F}'$ is the (ordinary) presheaf corresponding to
$\mathcal{F}$, the stalk $\mathcal{F}_x'$ is equal to
\[
  \varinjlim_{\mathfrak{B}}\mathcal{F}(V)
\]
over the set of the $V\in\mathfrak{B}$ containing $x$.
If $u:\mathcal{F}\to\mathcal{G}$ is a morphism of presheaves on
$\mathfrak{B}$ with values in $\mathbf{K}$, and
$u':\mathcal{F}'\to\mathcal{G}'$ the corresponding morphism of ordinary
presheaves, $u_x'$ is likewise the inductive limit of the morphisms
$u_V:\mathcal{F}(V)\to\mathcal{G}(V)$ for $V\in\mathfrak{B}$, $x\in V$.
\end{env}

\begin{env}[3.2.5]
\label{0.3.2.5}
Let us return to the general conditions of (3.2.1). If $\mathcal{F}$ is an
ordinary \emph{sheaf} with values in $\mathbf{K}$, and $\mathcal{F}_1$ the
sheaf \emph{on $\mathfrak{B}$} obtained by restricting $\mathcal{F}$ to
$\mathfrak{B}$, then the ordinary sheaf $\mathcal{F}_1'$ obtained from
$\mathcal{F}_1$ by the procedure of (3.2.1) is canonically isomorphic to
$\mathcal{F}$, by condition (F) and the uniqueness properties of the
projective limit. We shall usually identify $\mathcal{F}$ and
$\mathcal{F}_1'$.

If $\mathcal{G}$ is a second (ordinary) sheaf on $X$ with values in
$\mathbf{K}$, and $u:\mathcal{F}\to\mathcal{G}$ a morphism, the preceding
remark shows that the data of the
$u_V:\mathcal{F}(V)\to\mathcal{G}(V)$ \emph{for the $V\in\mathfrak{B}$
only} completely determine $u$; conversely, the $u_V$
being given for $V\in\mathfrak{B}$, it suffices to verify the commutativity diagram
with the restriction morphisms $\rho_V^W$ for $V\in\mathfrak{B}$,
$W\in\mathfrak{B}$ and $V\subset W$ in order that there exist one and only one
morphism $u'$ from $\mathcal{F}$ to $\mathcal{G}$ such that $u_V'=u_V$ for every
$V\in\mathfrak{B}$ (3.2.3).
\end{env}

\begin{env}[3.2.6]
\label{0.3.2.6}
Suppose again that $\mathbf{K}$ admits projective limits. Then the
category of \emph{sheaves on $X$ with values in $\mathbf{K}$} also admits
\emph{projective limits}; indeed, if $(\mathcal{F}_\lambda)$ is a
projective system of sheaves on $X$ with values in $\mathbf{K}$, the
\[
  \mathcal{F}(U)=\varprojlim_\lambda\mathcal{F}_\lambda(U)
\]
define a presheaf with values in $\mathbf{K}$, and the
verification of axiom (F) follows from the transitivity of projective
limits; the fact that $\mathcal{F}$ is then the projective limit of the
$\mathcal{F}_\lambda$ is immediate.

When $\mathbf{K}$ is the category of sets, for every
projective system $(\mathcal{H}_\lambda)$ such

\oldpage[0\textsubscript{I}]{28}

that $\mathcal{H}_\lambda$ is a \emph{subsheaf} of
$\mathcal{F}_\lambda$ for every $\lambda$,
$\varprojlim_\lambda\mathcal{H}_\lambda$ is canonically identified with a
\emph{subsheaf} of $\varprojlim_\lambda\mathcal{F}_\lambda$. If
$\mathbf{K}$ is the category of commutative groups, the covariant functor
$\varprojlim_\lambda\mathcal{F}_\lambda$ is \emph{additive and left exact}.
\end{env}

\subsection{Gluing sheaves}
\label{subsection:0.3.3}

\begin{env}[3.3.1]
\label{0.3.3.1}
Suppose again that the category $\mathbf{K}$ admits (generalized) projective limits.
Let $X$ be a topological space,
$\mathfrak{U}=(U_\lambda)_{\lambda\in L}$ an open covering of $X$, and
for each $\lambda\in L$, let $\mathcal{F}_\lambda$ be a sheaf on
$U_\lambda$ with values in $\mathbf{K}$; for every pair of indices
$(\lambda,\mu)$, suppose given an \emph{isomorphism}
\[
  \theta_{\lambda\mu}:
  \mathcal{F}_\mu|(U_\lambda\cap U_\mu)
  \xrightarrow{\sim}
  \mathcal{F}_\lambda|(U_\lambda\cap U_\mu);
\]
moreover, suppose that for every triple $(\lambda,\mu,\nu)$, denoting by
$\theta'_{\lambda\mu}$, $\theta'_{\mu\nu}$, $\theta'_{\lambda\nu}$ the
restrictions of $\theta_{\lambda\mu}$, $\theta_{\mu\nu}$,
$\theta_{\lambda\nu}$ to $U_\lambda\cap U_\mu\cap U_\nu$, we have
\[
  \theta'_{\lambda\nu}
  =\theta'_{\lambda\mu}\circ\theta'_{\mu\nu}
\]
(\emph{gluing condition for the $\theta_{\lambda\mu}$}). Then there
exist a sheaf $\mathcal{F}$ on $X$ with values in $\mathbf{K}$ and, for
each $\lambda$, an isomorphism
\[
  \eta_\lambda:\mathcal{F}|U_\lambda\xrightarrow{\sim}\mathcal{F}_\lambda
\]
such that, for every pair $(\lambda,\mu)$, denoting by
$\eta'_\lambda$ and $\eta'_\mu$ the restrictions of $\eta_\lambda$ and
$\eta_\mu$ to $U_\lambda\cap U_\mu$, we have
\[
  \theta_{\lambda\mu}=\eta'_\lambda\circ{\eta'_\mu}^{-1};
\]
moreover, $\mathcal{F}$ and the $\eta_\lambda$ are determined up to unique
isomorphism by these conditions. Indeed, the uniqueness follows
at once from (3.2.5). To establish the existence of $\mathcal{F}$, denote by
$\mathfrak{B}$ the basis of open subsets formed by the open subsets contained in at least one
$U_\lambda$, and for every $U\in\mathfrak{B}$ choose (by means of Hilbert's
$\tau$ function) one of the $\mathcal{F}_\lambda(U)$ for one of the
$\lambda$ such that $U\subset U_\lambda$; if this object is denoted by
$\mathcal{F}(U)$, the $\rho_U^V$ for $U\subset V$, $U\in\mathfrak{B}$,
$V\in\mathfrak{B}$ are defined in the obvious way (by means of the
$\theta_{\lambda\mu}$), and the transitivity condition is a consequence of the
gluing condition; moreover, the verification of
(F\textsubscript{0}) is immediate, so the presheaf on $\mathfrak{B}$
thus defined is indeed a sheaf, and from it one deduces by the general procedure
(3.2.1) an (ordinary) sheaf, again denoted $\mathcal{F}$, which answers the
question. We say that $\mathcal{F}$ is obtained by \emph{gluing the
$\mathcal{F}_\lambda$ by means of the $\theta_{\lambda\mu}$}, and we shall usually
identify $\mathcal{F}_\lambda$ and $\mathcal{F}|U_\lambda$ by means of
$\eta_\lambda$.

It is clear that every sheaf $\mathcal{F}$ on $X$ with values in
$\mathbf{K}$ can be regarded as obtained by gluing the sheaves
$\mathcal{F}_\lambda=\mathcal{F}|U_\lambda$ (where $(U_\lambda)$ is an
arbitrary open covering of $X$) by means of the isomorphisms
$\theta_{\lambda\mu}$ reduced to the identity.
\end{env}

\begin{env}[3.3.2]
\label{0.3.3.2}
With the same notation, let $\mathcal{G}_\lambda$ be a second sheaf on
$U_\lambda$ (for every $\lambda\in L$) with values in $\mathbf{K}$, and suppose
given, for every pair $(\lambda,\mu)$, an isomorphism
\[
  \omega_{\lambda\mu}:
  \mathcal{G}_\mu|(U_\lambda\cap U_\mu)
  \xrightarrow{\sim}
  \mathcal{G}_\lambda|(U_\lambda\cap U_\mu),
\]
these isomorphisms satisfying the gluing condition. Finally, suppose given,
for every $\lambda$, a morphism
$u_\lambda:\mathcal{F}_\lambda\to\mathcal{G}_\lambda$, and suppose that the diagrams
\begin{equation}
\begin{tikzcd}[column sep=huge,row sep=large]
\mathcal{F}_\mu|(U_\lambda\cap U_\mu)
  \arrow[r,"u_\mu"] \arrow[d]
  & \mathcal{G}_\mu|(U_\lambda\cap U_\mu) \arrow[d] \\
\mathcal{F}_\lambda|(U_\lambda\cap U_\mu)
  \arrow[r,"u_\lambda"']
  & \mathcal{G}_\lambda|(U_\lambda\cap U_\mu)
\end{tikzcd}
\tag{3.3.2.1}\label{0.3.3.2.1}
\end{equation}
are commutative. Then, if $\mathcal{G}$ is obtained by gluing the
$\mathcal{G}_\lambda$ by means of the $\omega_{\lambda\mu}$, there exists one and only one
morphism $u:\mathcal{F}\to\mathcal{G}$ such that the diagrams

\oldpage[0\textsubscript{I}]{29}

\[
\begin{tikzcd}[column sep=huge,row sep=large]
\mathcal{F}|U_\lambda
  \arrow[r,"u|U_\lambda"] \arrow[d]
  & \mathcal{G}|U_\lambda \arrow[d] \\
\mathcal{F}_\lambda
  \arrow[r,"u_\lambda"']
  & \mathcal{G}_\lambda
\end{tikzcd}
\]
are commutative; this follows at once from (3.2.3). The correspondence between
the family $(u_\lambda)$ and $u$ is a functorial bijection of the subset of
$\prod_\lambda\operatorname{Hom}(\mathcal{F}_\lambda,\mathcal{G}_\lambda)$
satisfying the conditions (3.3.2.1) onto
$\operatorname{Hom}(\mathcal{F},\mathcal{G})$.
\end{env}

\begin{env}[3.3.3]
\label{0.3.3.3}
With the notation of (3.3.1), let $V$ be an open subset of $X$; it is immediate that
the restrictions of the $\theta_{\lambda\mu}$ to $V\cap U_\lambda\cap U_\mu$
satisfy the gluing condition for the induced sheaves
$\mathcal{F}_\lambda|(V\cap U_\lambda)$, and that the sheaf on $V$ obtained by
gluing the latter is canonically identified with $\mathcal{F}|V$.
\end{env}

\subsection{Direct images of presheaves}
\label{subsection:0.3.4}

\begin{env}[3.4.1]
\label{0.3.4.1}
Let $X$, $Y$ be two topological spaces and $\psi:X\to Y$ a continuous
map. Let $\mathcal{F}$ be a presheaf on $X$ with values in a
category $\mathbf{K}$; for every open subset $U\subset Y$, let
\[
  \mathcal{G}(U)=\mathcal{F}(\psi^{-1}(U)),
\]
and if $U$, $V$ are two open subsets of $Y$ such that $U\subset V$, let
$\rho_U^V$ be the morphism
\[
  \mathcal{F}(\psi^{-1}(V))\to\mathcal{F}(\psi^{-1}(U));
\]
it is immediate that the $\mathcal{G}(U)$ and the $\rho_U^V$ define a
\emph{presheaf} on $Y$ with values in $\mathbf{K}$, which is called
the \emph{direct image of $\mathcal{F}$ by $\psi$} and which is denoted
$\psi_*(\mathcal{F})$. If $\mathcal{F}$ is a sheaf, one verifies at once
axiom (F) for the presheaf $\mathcal{G}(U)$, so
$\psi_*(\mathcal{F})$ is a sheaf.
\end{env}

\begin{env}[3.4.2]
\label{0.3.4.2}
Let $\mathcal{F}_1$, $\mathcal{F}_2$ be two presheaves on $X$ with values
in $\mathbf{K}$, and let $u:\mathcal{F}_1\to\mathcal{F}_2$ be a morphism.
As $U$ runs through the set of open subsets of $Y$, the family of
morphisms
\[
  u_{\psi^{-1}(U)}:
  \mathcal{F}_1(\psi^{-1}(U))\to\mathcal{F}_2(\psi^{-1}(U))
\]
satisfies the compatibility conditions with the restriction morphisms,
and consequently defines a morphism
\[
  \psi_*(u):\psi_*(\mathcal{F}_1)\to\psi_*(\mathcal{F}_2).
\]
If $v:\mathcal{F}_2\to\mathcal{F}_3$ is a morphism from $\mathcal{F}_2$ to
a third presheaf on $X$ with values in $\mathbf{K}$, we have
\[
  \psi_*(v\circ u)=\psi_*(v)\circ\psi_*(u);
\]
in other words, $\psi_*(\mathcal{F})$ is a \emph{covariant functor} in
$\mathcal{F}$, from the category of presheaves (resp. sheaves) on $X$ with
values in $\mathbf{K}$ to that of presheaves (resp. sheaves) on
$Y$ with values in $\mathbf{K}$.
\end{env}

\begin{env}[3.4.3]
\label{0.3.4.3}
Let $Z$ be a third topological space, $\psi':Y\to Z$ a continuous
map, and let $\psi''=\psi'\circ\psi$. It is clear that we have
\[
  \psi_*''(\mathcal{F})=\psi_*'(\psi_*(\mathcal{F}))
\]
for every presheaf $\mathcal{F}$ on $X$ with values in $\mathbf{K}$;
moreover, for every morphism $u:\mathcal{F}\to\mathcal{G}$ of such presheaves,
we have
\[
  \psi_*''(u)=\psi_*'(\psi_*(u)).
\]
In other words, $\psi_*''$ is the \emph{composite} of the functors $\psi_*'$
and $\psi_*$, which can be written
\[
  (\psi'\circ\psi)_*=\psi_*'\circ\psi_*.
\]

Moreover, for every open subset $U$ of $Y$, the image by the restriction
$\psi|\psi^{-1}(U)$ of the induced presheaf
$\mathcal{F}|\psi^{-1}(U)$ is none other than the induced presheaf
$\psi_*(\mathcal{F})|U$.
\end{env}

\begin{env}[3.4.4]
\label{0.3.4.4}
Suppose that the category $\mathbf{K}$ admits inductive limits, and let
$\mathcal{F}$ be a presheaf on $X$ with values in $\mathbf{K}$; for every
$x\in X$, the morphisms
\[
  \Gamma(\psi^{-1}(U),\mathcal{F})\to\mathcal{F}_x
\]
($U$ an open neighbourhood of $\psi(x)$ in $Y$) form an inductive system, which
by passage

\oldpage[0\textsubscript{I}]{30}

to the limit gives a morphism
\[
  \psi_x:(\psi_*(\mathcal{F}))_{\psi(x)}\to\mathcal{F}_x
\]
of stalks; in general, this morphism is \emph{neither injective nor surjective}.
It is functorial; indeed, if
$u:\mathcal{F}_1\to\mathcal{F}_2$ is a morphism of presheaves on $X$ with
values in $\mathbf{K}$, the diagram
\[
\begin{tikzcd}[column sep=huge,row sep=large]
(\psi_*(\mathcal{F}_1))_{\psi(x)}
  \arrow[r,"\psi_x"] \arrow[d,"{(\psi_*(u))_{\psi(x)}}"']
  & (\mathcal{F}_1)_x \arrow[d,"u_x"] \\
(\psi_*(\mathcal{F}_2))_{\psi(x)}
  \arrow[r,"\psi_x"']
  & (\mathcal{F}_2)_x
\end{tikzcd}
\]
is commutative. If $Z$ is a third topological space,
$\psi':Y\to Z$ a continuous map, and $\psi''=\psi'\circ\psi$, we have
\[
  \psi_x''=\psi_x\circ\psi'_{\psi(x)}
\]
for $x\in X$.
\end{env}

\begin{env}[3.4.5]
\label{0.3.4.5}
Under the hypotheses of (3.4.4), suppose moreover that $\psi$ is a
\emph{homeomorphism} of $X$ onto the subspace $\psi(X)$ of $Y$. Then, for
every $x\in X$, $\psi_x$ is an \emph{isomorphism}. This applies in
particular to the canonical injection $j$ of a subset $X$ of $Y$ into $Y$.
\end{env}

\begin{env}[3.4.6]
\label{0.3.4.6}
Suppose that $\mathbf{K}$ is the category of groups, or of rings, etc.
If $\mathcal{F}$ is a sheaf on $X$ with values in $\mathbf{K}$, with
support $S$, and if $y\notin\overline{\psi(S)}$, it follows from the definition of
$\psi_*(\mathcal{F})$ that
\[
  (\psi_*(\mathcal{F}))_y=\{0\},
\]
in other words, the support of $\psi_*(\mathcal{F})$ is contained in
$\overline{\psi(S)}$; but it is not necessarily contained in
$\psi(S)$. Under the same hypotheses, if $j$ is the canonical injection of a
subset $X$ of $Y$ into $Y$, the sheaf $j_*(\mathcal{F})$ induces
$\mathcal{F}$ on $X$; if moreover $X$ is \emph{closed} in $Y$,
$j_*(\mathcal{F})$ is the sheaf on $Y$ which induces $\mathcal{F}$ on $X$
and $0$ on $Y-X$ (G, II, 2.9.2), but in general it is distinct from the latter
when $X$ is assumed locally closed but not closed.
\end{env}

\subsection{Inverse images of presheaves}

\begin{env}[3.5.1]
\label{0.3.5.1}
Under the hypotheses of (3.4.1), if $\mathcal{F}$ (resp. $\mathcal{G}$) is a
presheaf on $X$ (resp. $Y$) with values in $\mathbf{K}$, every morphism
$u:\mathcal{G}\to\psi_*(\mathcal{F})$ of presheaves on $Y$ is also called
a $\psi$-\emph{morphism} from $\mathcal{G}$ to $\mathcal{F}$, and is
also written $\mathcal{G}\to\mathcal{F}$. We also denote by
$\operatorname{Hom}_\psi(\mathcal{G},\mathcal{F})$ the set
$\operatorname{Hom}_Y(\mathcal{G},\psi_*(\mathcal{F}))$ of the
$\psi$-morphisms from $\mathcal{G}$ to $\mathcal{F}$.

For every pair $(U,V)$, where $U$ is an open subset of $X$ and $V$ an open subset of $Y$
such that $\psi(U)\subset V$, one has a morphism
$u_{U,V}:\mathcal{G}(V)\to\mathcal{F}(U)$ obtained by composing the restriction
morphism
$\mathcal{F}(\psi^{-1}(V))\to\mathcal{F}(U)$ and the morphism
$u_V:\mathcal{G}(V)\to\psi_*(\mathcal{F})(V)=\mathcal{F}(\psi^{-1}(V))$; it
is immediate that these morphisms make commutative the diagrams
\begin{equation*}
\begin{tikzcd}[column sep=huge,row sep=large]
\mathcal{G}(V) \arrow[r,"u_{U,V}"] \arrow[d]
  & \mathcal{F}(U) \arrow[d] \\
\mathcal{G}(V') \arrow[r,"u_{U',V'}"']
  & \mathcal{F}(U')
\end{tikzcd}
\tag{3.5.1.1}\label{0.3.5.1.1}
\end{equation*}
for $U'\subset U$, $V'\subset V$, $\psi(U')\subset V'$. Conversely, a
family $(u_{U,V})$ of morphisms making the
diagrams (3.5.1.1) commutative defines a $\psi$-morphism $u$, since it suffices to take
\[
  u_V=u_{\psi^{-1}(V),V}.
\]

\oldpage[0\textsubscript{I}]{31}

If the category $\mathbf{K}$ admits (generalized) projective limits, and
if $\mathfrak{B}$, $\mathfrak{B}'$ are bases of the topologies of $X$ and
$Y$ respectively, then to define a $\psi$-morphism $u$ of \emph{sheaves},
one may restrict oneself to giving the $u_{U,V}$ for
$U\in\mathfrak{B}$, $V\in\mathfrak{B}'$ and $\psi(U)\subset V$, satisfying the
compatibility conditions (3.5.1.1) for $U$, $U'$ in $\mathfrak{B}$ and
$V$, $V'$ in $\mathfrak{B}'$; indeed, it suffices to define $u_W$, for
every open subset $W\subset Y$, as the projective limit of the $u_{U,V}$ for
$V\in\mathfrak{B}'$ and $V\subset W$, $U\in\mathfrak{B}$ and
$\psi(U)\subset V$.

When the category $\mathbf{K}$ admits inductive limits, one has, for
every $x\in X$, a morphism
\[
  \mathcal{G}(V)\to\mathcal{F}(\psi^{-1}(V))\to\mathcal{F}_x
\]
for every open neighbourhood $V$ of $\psi(x)$ in $Y$, and these morphisms
form an inductive system which by passage to the limit gives a morphism
\[
  \mathcal{G}_{\psi(x)}\to\mathcal{F}_x.
\]
\end{env}

\begin{env}[3.5.2]
\label{0.3.5.2}
Under the hypotheses of (3.4.3), let $\mathcal{F}$, $\mathcal{G}$,
$\mathcal{H}$ be presheaves with values in $\mathbf{K}$ on $X$, $Y$, $Z$
respectively, and let
$u:\mathcal{G}\to\psi_*(\mathcal{F})$,
$v:\mathcal{H}\to\psi'_*(\mathcal{G})$ be a $\psi$-morphism and a
$\psi'$-morphism respectively. One deduces from them a $\psi''$-morphism
\[
  w:\mathcal{H}\xrightarrow{v}\psi'_*(\mathcal{G})
  \xrightarrow{\psi'_*(u)}
  \psi'_*(\psi_*(\mathcal{F}))=\psi''_*(\mathcal{F}),
\]
which is called, by definition, the \emph{composite} of $u$ and $v$. One can
therefore regard the pairs $(X,\mathcal{F})$ consisting of a topological space
$X$ and a presheaf $\mathcal{F}$ on $X$ (with values in $\mathbf{K}$)
as forming a \emph{category}, the morphisms being the pairs
$(\psi,\theta):(X,\mathcal{F})\to(Y,\mathcal{G})$ consisting of a continuous
map $\psi:X\to Y$ and a $\psi$-morphism
$\theta:\mathcal{G}\to\mathcal{F}$.
\end{env}

\begin{env}[3.5.3]
\label{0.3.5.3}
Let $\psi:X\to Y$ be a continuous map and $\mathcal{G}$ a
\emph{presheaf} on $Y$ with values in $\mathbf{K}$. We shall call an
\emph{inverse image of $\mathcal{G}$ by $\psi$} a pair
$(\mathcal{G}',\rho)$, where $\mathcal{G}'$ is a \emph{sheaf} on $X$ with
values in $\mathbf{K}$ and
$\rho:\mathcal{G}\to\mathcal{G}'$ a $\psi$-morphism (in other words a
homomorphism $\mathcal{G}\to\psi_*(\mathcal{G}'))$ such that, for every
\emph{sheaf} $\mathcal{F}$ on $X$ with values in $\mathbf{K}$,
the map
\begin{equation*}
  \operatorname{Hom}_X(\mathcal{G}',\mathcal{F})
  \longrightarrow \operatorname{Hom}_\psi(\mathcal{G},\mathcal{F})
  =\operatorname{Hom}_Y(\mathcal{G},\psi_*(\mathcal{F}))
\tag{3.5.3.1}\label{0.3.5.3.1}
\end{equation*}
sending $v$ to $\psi_*(v)\circ\rho$ is a \emph{bijection}; this
map, being functorial in $\mathcal{F}$, will then define an
isomorphism of functors in $\mathcal{F}$. Since the pair
$(\mathcal{G}',\rho)$ is a solution of a universal problem, we know that it
is \emph{determined up to unique isomorphism} when it exists. We shall
then write $\mathcal{G}'=\psi^*(\mathcal{G})$, $\rho=\rho_{\mathcal{G}}$, and, by
abuse of language, we shall say that $\psi^*(\mathcal{G})$ is \emph{the inverse
image sheaf} of $\mathcal{G}$ by $\psi$, it being understood that
$\psi^*(\mathcal{G})$ is regarded as equipped with the \emph{canonical}
$\psi$-morphism
$\rho_{\mathcal{G}}:\mathcal{G}\to\psi^*(\mathcal{G})$, that is, with
the \emph{canonical homomorphism} of presheaves on $Y$:
\begin{equation*}
  \rho_{\mathcal{G}}:\mathcal{G}\to\psi_*(\psi^*(\mathcal{G})).
\tag{3.5.3.2}\label{0.3.5.3.2}
\end{equation*}
For every homomorphism
$v:\psi^*(\mathcal{G})\to\mathcal{F}$ (where $\mathcal{F}$ is a sheaf on
$X$ with values in $\mathbf{K}$), we shall set
$v^b=\psi_*(v)\circ\rho_{\mathcal{G}}:
\mathcal{G}\to\psi_*(\mathcal{F})$. By definition, \emph{every} morphism of
presheaves $u:\mathcal{G}\to\psi_*(\mathcal{F})$ is of the form $v^b$ for
one and only one $v$, which we shall denote by $u^\sharp$. In other words, every
morphism $u:\mathcal{G}\to\psi_*(\mathcal{F})$ of presheaves factorizes
uniquely as
\begin{equation*}
  u:\mathcal{G}\xrightarrow{\rho_{\mathcal{G}}}
  \psi_*(\psi^*(\mathcal{G}))
  \xrightarrow{\psi_*(u^\sharp)}\psi_*(\mathcal{F}).
\tag{3.5.3.3}\label{0.3.5.3.3}
\end{equation*}
\end{env}

\oldpage[0\textsubscript{I}]{32}

\begin{env}[3.5.4]
\label{0.3.5.4}
Suppose now that the category $\mathbf{K}$ is such\footnote{In
the book cited in the Introduction, we shall give very
general conditions on the category $\mathbf{K}$ ensuring the existence of inverse
images of presheaves with values in $\mathbf{K}$.} that \emph{every}
presheaf $\mathcal{G}$ on $Y$ with values in $\mathbf{K}$ admits an
inverse image by $\psi$, which we shall denote by $\psi^*(\mathcal{G})$.

We shall see that one can define $\psi^*(\mathcal{G})$ as a \emph{covariant
functor} in $\mathcal{G}$, from the category of presheaves on $Y$ with
values in $\mathbf{K}$ to that of sheaves on $X$ with values in
$\mathbf{K}$, in such a way that the isomorphism $v\to v^b$ is an
\emph{isomorphism of bifunctors}
\begin{equation*}
  \operatorname{Hom}_X(\psi^*(\mathcal{G}),\mathcal{F})
  \xrightarrow{\sim}
  \operatorname{Hom}_Y(\mathcal{G},\psi_*(\mathcal{F}))
\tag{3.5.4.1}\label{0.3.5.4.1}
\end{equation*}
in $\mathcal{G}$ and $\mathcal{F}$.

Indeed, for every morphism
$w:\mathcal{G}_1\to\mathcal{G}_2$ of presheaves on $Y$ with values in
$\mathbf{K}$, consider the composite morphism
\[
  \mathcal{G}_1\xrightarrow{w}\mathcal{G}_2
  \xrightarrow{\rho_{\mathcal{G}_2}}
  \psi_*(\psi^*(\mathcal{G}_2))~;
\]
there corresponds to it a morphism
$(\rho_{\mathcal{G}_2}\circ w)^\sharp:
\psi^*(\mathcal{G}_1)\to\psi^*(\mathcal{G}_2)$, which we shall denote by
$\psi^*(w)$. We therefore have, by (3.5.3.3),
\begin{equation*}
  \psi_*(\psi^*(w))\circ\rho_{\mathcal{G}_1}
  =\rho_{\mathcal{G}_2}\circ w.
\tag{3.5.4.2}\label{0.3.5.4.2}
\end{equation*}
For every morphism $u:\mathcal{G}_2\to\psi_*(\mathcal{F})$, where
$\mathcal{F}$ is a sheaf on $X$ with values in $\mathbf{K}$, we have,
by (3.5.3.3), (3.5.4.2) and the definition of $v^b$,
\[
  (u^\sharp\circ\psi^*(w))^b
  =\psi_*(u^\sharp)\circ\psi_*(\psi^*(w))\circ\rho_{\mathcal{G}_1}
  =\psi_*(u^\sharp)\circ\rho_{\mathcal{G}_2}\circ w
  =u\circ w
\]
or again
\begin{equation*}
  (u\circ w)^\sharp=u^\sharp\circ\psi^*(w).
\tag{3.5.4.3}\label{0.3.5.4.3}
\end{equation*}
If in particular we take for $u$ a morphism
\[
  \mathcal{G}_2\xrightarrow{w'}\mathcal{G}_3
  \xrightarrow{\rho_{\mathcal{G}_3}}
  \psi_*(\psi^*(\mathcal{G}_3)),
\]
we obtain
\[
  \psi^*(w'\circ w)
  =(\rho_{\mathcal{G}_3}\circ w'\circ w)^\sharp
  =(\rho_{\mathcal{G}_3}\circ w')^\sharp\circ\psi^*(w)
  =\psi^*(w')\circ\psi^*(w),
\]
whence our assertion.

Finally, for every sheaf $\mathcal{F}$ on $X$ with values in
$\mathbf{K}$, let $i_{\mathcal{F}}$ be the identity morphism of
$\psi_*(\mathcal{F})$ and denote by
\[
  \sigma_{\mathcal{F}}:
  \psi^*(\psi_*(\mathcal{F}))\to\mathcal{F}
\]
the morphism $(i_{\mathcal{F}})^\sharp$; formula (3.5.4.3) gives in
particular the factorization
\begin{equation*}
  u^\sharp:\psi^*(\mathcal{G})
  \xrightarrow{\psi^*(u)}
  \psi^*(\psi_*(\mathcal{F}))
  \xrightarrow{\sigma_{\mathcal{F}}}\mathcal{F}
\tag{3.5.4.4}\label{0.3.5.4.4}
\end{equation*}
for every morphism $u:\mathcal{G}\to\psi_*(\mathcal{F})$. We shall say that the
morphism $\sigma_{\mathcal{F}}$ is \emph{canonical}.
\end{env}

\begin{env}[3.5.5]
\label{0.3.5.5}
Let $\psi':Y\to Z$ be a continuous map, and suppose that every
presheaf $\mathcal{H}$ on $Z$ with values in $\mathbf{K}$ admits an
inverse image $\psi'^*(\mathcal{H})$ by $\psi'$. Then (with the
hypotheses of (3.5.4)) every presheaf $\mathcal{H}$ on $Z$ with values in
$\mathbf{K}$ admits an inverse image by $\psi''=\psi'\circ\psi$, and there is
a canonical functorial isomorphism
\begin{equation*}
  \psi''^*(\mathcal{H})\xrightarrow{\sim}
  \psi^*(\psi'^*(\mathcal{H}))
\tag{3.5.5.1}\label{0.3.5.5.1}
\end{equation*}

\oldpage[0\textsubscript{I}]{33}

This indeed follows at once from the definitions, taking into account that
$\psi''_*=\psi'_*\circ\psi_*$. Moreover, if
$u:\mathcal{G}\to\psi_*(\mathcal{F})$ is a $\psi$-morphism,
$v:\mathcal{H}\to\psi'_*(\mathcal{G})$ a $\psi'$-morphism, and
$w=\psi'_*(u)\circ v$ their composite (3.5.2), one sees at once that
$w^\sharp$ is the composite morphism
\[
  w^\sharp:\psi^*(\psi'^*(\mathcal{H}))
  \xrightarrow{\psi^*(v^\sharp)}
  \psi^*(\mathcal{G})
  \xrightarrow{u^\sharp}\mathcal{F}.
\]
\end{env}

\begin{env}[3.5.6]
\label{0.3.5.6}
Take in particular for $\psi$ the identity map
$1_X:X\to X$. Then, if the inverse image by $\psi$ of a presheaf
$\mathcal{F}$ on $X$ with values in $\mathbf{K}$ exists, this
inverse image is called the \emph{sheaf associated with the presheaf $\mathcal{F}$}.
Every morphism $u:\mathcal{F}\to\mathcal{F}'$ from $\mathcal{F}$ to a
\emph{sheaf} $\mathcal{F}'$ with values in $\mathbf{K}$ therefore factorizes
uniquely as
\[
  \mathcal{F}\xrightarrow{\rho_{\mathcal{F}}}
  1_X^*(\mathcal{F})\xrightarrow{u^\sharp}\mathcal{F}'.
\]
\end{env}


\subsection{Simple sheaves and locally simple sheaves}

\begin{env}[3.6.1]
\label{0.3.6.1}
We shall say that a \emph{presheaf} $\mathcal{F}$ on $X$ with values in
$\mathbf{K}$ is \emph{constant} if the canonical morphisms
$\mathcal{F}(X)\to\mathcal{F}(U)$ are \emph{isomorphisms} for every
non-empty open set $U\subset X$; note that $\mathcal{F}$ is not
necessarily a sheaf. A \emph{sheaf} is said to be \emph{simple}
if it is associated (3.5.6) to a constant presheaf. A sheaf
$\mathcal{F}$ is said to be \emph{locally simple} if every $x\in X$ has an
open neighbourhood $U$ such that $\mathcal{F}|U$ is simple.
\end{env}

\begin{env}[3.6.2]
\label{0.3.6.2}
Suppose that $X$ is \emph{irreducible} (2.1.1); then the following
properties are equivalent:
\begin{enumerate}
  \item[a)] \emph{$\mathcal{F}$ is a constant presheaf on $X$};
  \item[b)] \emph{$\mathcal{F}$ is a simple sheaf on $X$};
  \item[c)] \emph{$\mathcal{F}$ is a locally simple sheaf on $X$}.
\end{enumerate}

Indeed, let $\mathcal{F}$ be a constant presheaf on $X$; if $U$, $V$
are two non-empty open sets in $X$, then $U\cap V$ is non-empty, so, since
$\mathcal{F}(X)\to\mathcal{F}(U)\to\mathcal{F}(U\cap V)$ and
$\mathcal{F}(X)\to\mathcal{F}(U)$ are isomorphisms, so is
$\mathcal{F}(U)\to\mathcal{F}(U\cap V)$, and likewise
$\mathcal{F}(V)\to\mathcal{F}(U\cap V)$ is an isomorphism. One concludes
at once that axiom (F) of (3.1.2) is indeed satisfied, that $\mathcal{F}$ is
isomorphic to its associated sheaf, and consequently that a) implies b).

Now let $(U_\alpha)$ be an open covering of $X$ by non-empty open sets,
and let $\mathcal{F}$ be a sheaf on $X$ such that
$\mathcal{F}|U_\alpha$ is simple for every $\alpha$; since $U_\alpha$ is
irreducible, $\mathcal{F}|U_\alpha$ is a constant presheaf by
the foregoing. Since $U_\alpha\cap U_\beta$ is not empty,
$\mathcal{F}(U_\alpha)\to\mathcal{F}(U_\alpha\cap U_\beta)$ and
$\mathcal{F}(U_\beta)\to\mathcal{F}(U_\alpha\cap U_\beta)$ are
isomorphisms, from which one deduces a canonical isomorphism
$\theta_{\alpha\beta}:\mathcal{F}(U_\alpha)\to\mathcal{F}(U_\beta)$ for every
pair of indices. But then, if one applies condition (F) for $U=X$, one
sees that for every index $\alpha_0$, $\mathcal{F}(U_{\alpha_0})$ and the
$\theta_{\alpha_0\alpha}$ form a solution of the universal problem, which
(by uniqueness) implies that
$\mathcal{F}(X)\to\mathcal{F}(U_{\alpha_0})$ is an isomorphism, and therefore
proves that c) implies a).
\end{env}

\subsection{Inverse images of presheaves of groups or of rings}

\oldpage[0\textsubscript{I}]{34}

\begin{env}[3.7.1]
\label{0.3.7.1}
We shall show that when $\mathbf{K}$ is taken to be the category of
sets, the inverse image under $\psi$ of every presheaf $\mathcal{G}$ with
values in $\mathbf{K}$ \emph{always exists} (the notation and hypotheses
on $X$, $Y$, $\psi$ being those of (3.5.3)). Indeed, for every open set
$U\subset X$, define $\mathcal{G}'(U)$ as follows: an element $s'$ of
$\mathcal{G}'(U)$ is a family $(s'_x)_{x\in U}$, where
$s'_x\in\mathcal{G}_{\psi(x)}$ for every $x\in U$, and where, for every $x\in U$,
the following condition is satisfied: there exist an open neighbourhood $V$ of
$\psi(x)$ in $Y$, a neighbourhood $W\subset\psi^{-1}(V)\cap U$ of $x$ and an
element $s\in\mathcal{G}(V)$ such that $s'_z=s_{\psi(z)}$ for every $z\in W$.
One verifies immediately that $U\mapsto\mathcal{G}'(U)$ does satisfy the
axioms of \emph{sheaves}.

Now let $\mathcal{F}$ be a sheaf of sets on $X$, and let
$u:\mathcal{G}\to\psi_*(\mathcal{F})$,
$v:\mathcal{G}'\to\mathcal{F}$ be morphisms. One defines $u^\sharp$ and $v^b$
as follows: if $s'$ is a section of $\mathcal{G}'$ over
a neighbourhood $U$ of $x\in X$, and if $V$ is an open neighbourhood of $\psi(x)$
and $s\in\mathcal{G}(V)$ is such that $s'_z=s_{\psi(z)}$ for $z$ in a
neighbourhood of $x$ contained in $\psi^{-1}(V)\cap U$, one takes
$u_x^\sharp(s'_x)=u_{\psi(x)}(s_{\psi(x)})$. Likewise, if
$s\in\mathcal{G}(V)$ ($V$ open in $Y$), $v^b(s)$ is the section of
$\mathcal{F}$ over $\psi^{-1}(V)$ that is the image under $v$ of the section $s'$ of
$\mathcal{G}'$ such that $s'_x=s_{\psi(x)}$ for every
$x\in\psi^{-1}(V)$. Moreover, the canonical homomorphism (3.5.3)
$\rho:\mathcal{G}\to\psi_*(\psi^*(\mathcal{G}))$ is defined as
follows: for every open set $V\subset Y$ and every section
$s\in\Gamma(V,\mathcal{G})$, $\rho(s)$ is the section
$(s_{\psi(x)})_{x\in\psi^{-1}(V)}$ of $\psi^*(\mathcal{G})$ over
$\psi^{-1}(V)$. The verification of the relations
$(u^\sharp)^b=u$, $(v^b)^\sharp=v$ and
$v^b=\psi_*(v)\circ\rho$ is immediate, and proves our assertion.

One verifies that, if $w:\mathcal{G}_1\to\mathcal{G}_2$ is a homomorphism of
presheaves of sets on $Y$, then $\psi^*(w)$ is given explicitly as
follows: if $s'=(s'_x)_{x\in U}$ is a section of
$\psi^*(\mathcal{G}_1)$ over an open set $U$ of $X$,
$(\psi^*(w))(s')$ is the family
$(w_{\psi(x)}(s'_x))_{x\in U}$. Finally, it is immediate that for every open set
$V$ of $Y$, the inverse image of $\mathcal{G}|V$ under the restriction of
$\psi$ to $\psi^{-1}(V)$ is identical to the induced sheaf
$\psi^*(\mathcal{G})|\psi^{-1}(V)$.

When $\psi$ is the identity $1_X$, one recovers the definition of a sheaf
of sets associated to a presheaf (G, II, 1.2). The preceding
considerations apply without change when $\mathbf{K}$ is the
category of groups or of rings (not necessarily commutative).

When $X$ is an arbitrary subset of a topological space $Y$, and $j$ is
the canonical injection $X\to Y$, then for every sheaf $\mathcal{G}$ on $Y$ with
values in a category $\mathbf{K}$, the inverse image $j^*(\mathcal{G})$ (when it
exists) is called the sheaf \emph{induced} on $X$ by $\mathcal{G}$; for sheaves
of sets (or of groups, or of rings) one recovers the usual definition
(G, II, 1.5).
\end{env}

\begin{env}[3.7.2]
\label{0.3.7.2}
Keeping the notation and hypotheses of (3.5.3), suppose that
$\mathcal{G}$ is a \emph{sheaf} of groups (resp. of rings) on $Y$.
The definition of the sections of $\psi^*(\mathcal{G})$ (3.7.1) shows (taking
(3.4.4) into account) that the homomorphism of stalks
$\psi_x\circ\rho_{\psi(x)}:\mathcal{G}_{\psi(x)}\to
(\psi^*(\mathcal{G}))_x$ is an \emph{isomorphism functorial in
$\mathcal{G}$}, which allows one to identify these two stalks; with this
identification, $u_x^\sharp$ is identical to the homomorphism defined in
(3.5.1), and in particular one has
$\operatorname{Supp}(\psi^*(\mathcal{G}))=
\psi^{-1}(\operatorname{Supp}(\mathcal{G}))$.

An immediate consequence of this result is that \emph{the functor
$\psi^*(\mathcal{G})$ is exact in $\mathcal{G}$} in the abelian category
of sheaves of commutative groups.
\end{env}

\subsection{Sheaves of pseudo-discrete spaces}

\oldpage[0\textsubscript{I}]{35}

\begin{env}[3.8.1]
\label{0.3.8.1}
Let $X$ be a topological space whose topology has a base
$\mathfrak{B}$ consisting of \emph{quasi-compact} open sets. Let
$\mathcal{F}$ be a \emph{sheaf of sets} on $X$; if each of the
$\mathcal{F}(U)$ is endowed with the \emph{discrete} topology,
$U\mapsto\mathcal{F}(U)$ is a \emph{presheaf of topological spaces}.
We shall see that there exists a \emph{sheaf of topological spaces
$\mathcal{F}'$ associated to $\mathcal{F}$} (3.5.6) such that
$\Gamma(U,\mathcal{F}')$ is the discrete space $\mathcal{F}(U)$ for every
\emph{quasi-compact} open set $U$. For this it will suffice to show that the
presheaf $U\mapsto\mathcal{F}(U)$ of discrete topological spaces
\emph{on $\mathfrak{B}$} satisfies condition $(\mathrm{F}_0)$ of (3.2.2),
and more generally that if $U$ is a quasi-compact open set and
$(U_\alpha)$ is a covering of $U$ by sets of $\mathfrak{B}$,
the coarsest topology $\mathcal{T}$ on $\Gamma(U,\mathcal{F})$ making
the maps
$\Gamma(U,\mathcal{F})\to\Gamma(U_\alpha,\mathcal{F})$ continuous is the
\emph{discrete} topology. Now, there exist finitely many indices $\alpha_i$ such that
$U=\bigcup_i U_{\alpha_i}$. Let $s\in\Gamma(U,\mathcal{F})$ and let $s_i$ be its
image in $\Gamma(U_{\alpha_i},\mathcal{F})$; the intersection of the inverse
images of the sets $\{s_i\}$ is by definition a neighbourhood of $s$
for $\mathcal{T}$; but since $\mathcal{F}$ is a sheaf of sets and
the $U_{\alpha_i}$ cover $U$, this intersection reduces to $s$,
whence our assertion.

Note that if $U$ is a non-quasi-compact open set of $X$, the topological
space $\Gamma(U,\mathcal{F}')$ still has $\Gamma(U,\mathcal{F})$ as its
underlying set, but its topology is not discrete in general:
it is the coarsest topology making the maps
$\Gamma(U,\mathcal{F})\to\Gamma(V,\mathcal{F})$ continuous, for $V\in\mathfrak{B}$ and
$V\subset U$ (the $\Gamma(V,\mathcal{F})$ being discrete).

The preceding considerations apply without modification to sheaves
of groups or of rings (not necessarily commutative), and associate with them
sheaves of \emph{topological groups} or of \emph{topological rings},
respectively. For brevity, we shall say that the sheaf $\mathcal{F}'$ is the
sheaf of \emph{pseudo-discrete} \emph{spaces} (resp. \emph{groups},
\emph{rings}) associated to the sheaf of sets (resp. groups,
rings) $\mathcal{F}$.
\end{env}

\begin{env}[3.8.2]
\label{0.3.8.2}
Let $\mathcal{F}$, $\mathcal{G}$ be two sheaves of sets (resp. groups,
rings) on $X$, and $u:\mathcal{F}\to\mathcal{G}$ a homomorphism. Then $u$
is also a \emph{continuous} homomorphism $\mathcal{F}'\to\mathcal{G}'$, where
$\mathcal{F}'$ and $\mathcal{G}'$ denote the pseudo-discrete sheaves
associated to $\mathcal{F}$ and $\mathcal{G}$; indeed, this follows from
(3.2.5).
\end{env}

\begin{env}[3.8.3]
\label{0.3.8.3}
Let $\mathcal{F}$ be a sheaf of sets, $\mathcal{H}$ a sub-sheaf of
$\mathcal{F}$, and $\mathcal{F}'$ and $\mathcal{H}'$ the pseudo-discrete sheaves
associated to $\mathcal{F}$ and $\mathcal{H}$ respectively. Then, for every
open set $U\subset X$, $\Gamma(U,\mathcal{H}')$ is \emph{closed} in
$\Gamma(U,\mathcal{F}')$: indeed, it is the intersection of the inverse
images of the $\Gamma(V,\mathcal{H})$ (for $V\in\mathfrak{B}$,
$V\subset U$) under the continuous maps
$\Gamma(U,\mathcal{F})\to\Gamma(V,\mathcal{F})$, and
$\Gamma(V,\mathcal{H})$ is closed in the discrete space
$\Gamma(V,\mathcal{F})$.
\end{env}

